\Askhsh{36097}{2}
\begin{ylh}{1}
    \item 3.1. Είδη και στοιχεία τριγώνων
    \item 4.2. Τέμνουσα δύο ευθειών - Ευκλείδειο αίτημα
    \item 4.6. Άθροισμα γωνιών τριγώνου
    \item 5.9. Μια ιδιότητα του ορθογώνιου τριγώνου.
\end{ylh}
% ===================================================================
%  Αρχείο: 36097.tex (Fragment)
%  Αυτόματη μετατροπή μέσω doc2latex.py
% ===================================================================

ΘΕΜΑ 2

Στο παρακάτω σχήμα οι ευθείες ε\textsubscript{1}, ε\textsubscript{2} είναι παράλληλες και AB=6.

\begin{alist}
\item Να υπολογίσετε τις γωνίες $\widehat{\varphi}$ και $\widehat{\omega}$. \monades{10}

\item Να προσδιορίσετε το είδος του τριγώνου ΑΒΚ ως προς τις γωνίες του. \monades{7}

\item Να υπολογίσετε το μήκος της ΑΚ, αιτιολογώντας την απάντησή σας. \monades{8}


\begin{center}
\begin{tikzpicture}[scale=0.8]
\tkzDefPoint(0,2){k1}
\tkzDefPoint(6,2){k2}
\tkzDefPoint(0,0){l1}
\tkzDefPoint(6,0){l2}

\tkzDefPoint(1,2){A}
\tkzDefPoint(4,2){K}
\tkzDefPoint(2.5,0){B}

\draw[l3] (k1)--(k2) node[right]{$\epsilon_1$};
\draw[l3] (l1)--(l2) node[right]{$\epsilon_2$};
\draw[l3] (A)--(B)--(K)--cycle;

\tkzDrawPoints(A,B,K)
\tkzLabelPoint[above](A){$A$}
\tkzLabelPoint[above](K){$K$}
\tkzLabelPoint[below](B){$B$}

\tkzMarkAngle[size=0.6,mark=none](B,A,K)
\tkzLabelAngle[pos=0.8](B,A,K){$\varphi$}
\tkzMarkAngle[size=0.6,mark=none](K,B,A) % Assuming omega here
\tkzMarkAngle[size=0.6,mark=none](A,B,l1)
\tkzLabelAngle[pos=0.8](A,B,l1){$\omega$}

\tkzLabelSegment[above](A,B){$6$}
\end{tikzpicture}
\end{center}
\end{alist}
